% p201_1; source PDF page 201; manually verified editable table.
\begin{tikzpicture}[x=1cm,y=-1cm,inner sep=0pt,outer sep=0pt]
\path[use as bounding box] (0,0) rectangle (16.000,5.900);
\draw[line width=.4pt] (0.000,0.000) rectangle (4.500,0.850);
\node[align=center,text width=4.240cm,font={\linespread{1}\fontsize{10}{13.50}\selectfont\bfseries}] at (2.250,0.425) {种类};
\draw[line width=.4pt] (4.500,0.000) rectangle (8.500,0.850);
\node[align=center,text width=3.740cm,font={\linespread{1}\fontsize{10}{13.50}\selectfont\bfseries}] at (6.500,0.425) {投案方式};
\draw[line width=.4pt] (8.500,0.000) rectangle (16.000,0.850);
\node[align=center,text width=7.240cm,font={\linespread{1}\fontsize{10}{13.50}\selectfont\bfseries}] at (12.250,0.425) {供述内容};
\draw[line width=.4pt] (0.000,0.850) rectangle (4.500,1.800);
\node[align=center,text width=4.240cm,font={\linespread{1}\fontsize{10}{13.50}\selectfont}] at (2.250,1.325) {一般自首};
\draw[line width=.4pt] (4.500,0.850) rectangle (8.500,1.800);
\node[align=center,text width=3.740cm,font={\linespread{1}\fontsize{10}{13.50}\selectfont}] at (6.500,1.325) {因A罪行自动投案};
\draw[line width=.4pt] (8.500,0.850) rectangle (16.000,1.800);
\node[align=left,text width=7.240cm,font={\linespread{1}\fontsize{10}{13.50}\selectfont}] at (12.250,1.325) {如实供述A罪行（司法机关可掌握、可未掌握）};
\draw[line width=.4pt] (0.000,1.800) rectangle (4.500,2.750);
\node[align=center,text width=4.240cm,font={\linespread{1}\fontsize{10}{13.50}\selectfont}] at (2.250,2.275) {特别自首};
\draw[line width=.4pt] (4.500,1.800) rectangle (8.500,2.750);
\node[align=center,text width=3.740cm,font={\linespread{1}\fontsize{10}{13.50}\selectfont}] at (6.500,2.275) {因A罪行被动归案};
\draw[line width=.4pt] (8.500,1.800) rectangle (16.000,2.750);
\node[align=left,text width=7.240cm,font={\linespread{1}\fontsize{10}{13.50}\selectfont}] at (12.250,2.275) {如实供述司法机关尚未掌握的B罪行};
\draw[line width=.4pt] (0.000,2.750) rectangle (4.500,3.700);
\node[align=center,text width=4.240cm,font={\linespread{1}\fontsize{10}{13.50}\selectfont}] at (2.250,3.225) {坦白};
\draw[line width=.4pt] (4.500,2.750) rectangle (8.500,3.700);
\node[align=center,text width=3.740cm,font={\linespread{1}\fontsize{10}{13.50}\selectfont}] at (6.500,3.225) {因A罪行被动归案};
\draw[line width=.4pt] (8.500,2.750) rectangle (16.000,3.700);
\node[align=left,text width=7.240cm,font={\linespread{1}\fontsize{10}{13.50}\selectfont}] at (12.250,3.225) {如实供述司法机关已经掌握的A罪行};
\draw[line width=.4pt] (0.000,3.700) rectangle (2.500,5.900);
\node[align=center,text width=2.240cm,font={\linespread{1}\fontsize{10}{13.50}\selectfont}] at (1.250,4.800) {甲乙构成A罪的共同犯罪};
\draw[line width=.4pt] (2.500,3.700) rectangle (4.500,4.800);
\node[align=center,text width=1.740cm,font={\linespread{1}\fontsize{10}{13.50}\selectfont}] at (3.500,4.250) {一般自首};
\draw[line width=.4pt] (4.500,3.700) rectangle (8.500,4.800);
\node[align=center,text width=3.740cm,font={\linespread{1}\fontsize{10}{13.50}\selectfont}] at (6.500,4.250) {甲自动投案};
\draw[line width=.4pt] (8.500,3.700) rectangle (16.000,4.800);
\node[align=center,text width=7.240cm,font={\linespread{1}\fontsize{10}{13.50}\selectfont}] at (12.250,4.250) {如实供述乙是同案犯};
\draw[line width=.4pt] (2.500,4.800) rectangle (4.500,5.900);
\node[align=center,text width=1.740cm,font={\linespread{1}\fontsize{10}{13.50}\selectfont}] at (3.500,5.350) {立功};
\draw[line width=.4pt] (4.500,4.800) rectangle (8.500,5.900);
\node[align=center,text width=3.740cm,font={\linespread{1}\fontsize{10}{13.50}\selectfont}] at (6.500,5.350) {甲自动投案或被动归案};
\draw[line width=.4pt] (8.500,4.800) rectangle (16.000,5.900);
\node[align=left,text width=7.240cm,font={\linespread{1}\fontsize{10}{13.50}\selectfont}] at (12.250,5.350) {揭发司法机关尚未掌握的乙的B罪行};
\end{tikzpicture}
